\begin{table}[htbp]
\centering
\caption{Opening-price digit, volatility, and measured clustering}
\label{tab:opening}
\small
\begin{threeparttable}
\begin{tabular}{@{}lrrrr@{}}
\toprule
Opening digit & Low volatility (\%) & High volatility (\%) & Difference & Stock-days, low vol. \\
\midrule
0 & 14.37 & 15.95 & -1.58 & 72 \\
\bottomrule
\end{tabular}
\begin{tablenotes}[flushleft]\footnotesize
\item Mean $M^{0}$ by the last digit of the opening price, split on whether the previous day's realised variance was below that day's cross-sectional median. On quiet days a stock that opens far from a round price may never reach one, which depresses the measure mechanically. Every regression in this study therefore controls for the interaction of the opening-digit dummies with the low-volatility indicator, as Ohta (2026) does.
\end{tablenotes}
\end{threeparttable}
\end{table}
